%% ma: L12-B01-0040
%% lop: 12
%% chuong: 1
%% bai: 1
%% dang: 12.1.9
%% loai: TN4
%% muc: TH
%% dap_an: "B"
%% nguon: "Ảnh giáo viên gửi 05/11/2026 (M11): Kiểm tra giữa HK1 Toán 12 (2026-2027), THPT Hồ Thị Bi (TP.HCM), mã đề 1133, Phần I Câu 6"
%% kien_thuc: "Bài 1 (điểm cực trị: dấu của đạo hàm, nghiệm bội chẵn không đổi dấu)"
%% trang_thai: chinh-thuc
%% ghi_chu: "Đề gốc không kèm đáp án; đã giải độc lập. kho.py trung gợi ý thay số so với L12-B01-0022 (khung 0,81) nhưng khác nội dung (đạo hàm có nghiệm bội chẵn); giữ lại."
\begin{ex}
Cho hàm số $y=f(x)$ có đạo hàm $f'(x)=x(x-1)^2(x+2)^4$ với mọi $x\in\mathbb R$. Số điểm cực trị của hàm số đã cho là
\choice{$2$}{\True $1$}{$3$}{$0$}
\loigiai{$f'(x)=0\Leftrightarrow x=0$, $x=1$, $x=-2$.
$x=1$ và $x=-2$ là nghiệm bội chẵn nên $f'(x)$ không đổi dấu khi qua các điểm đó.
$x=0$ là nghiệm đơn, $f'(x)$ đổi dấu từ âm sang dương. Vậy hàm số có đúng $1$ điểm cực trị.}
\end{ex}
